% Guion 20 min — Efecto Mpemba cuántico en HPC. Compilar: pdflatex guion_20min
\documentclass[11pt,a4paper]{article}
\usepackage[utf8]{inputenc}
\usepackage[T1]{fontenc}
\usepackage[spanish,es-noshorthands]{babel}
\usepackage[margin=2.2cm]{geometry}
\usepackage{xcolor}
\usepackage{enumitem}
\usepackage{amssymb}
\usepackage{hyperref}
\hypersetup{hidelinks}

\definecolor{juan}{RGB}{15,115,140}
\definecolor{seb}{RGB}{205,95,25}
\definecolor{ink}{RGB}{33,37,48}
\definecolor{soft}{RGB}{120,125,135}
\color{ink}

\newcommand{\Jb}{\colorbox{juan}{\color{white}\bfseries\footnotesize~Juan~}}
\newcommand{\Sb}{\colorbox{seb}{\color{white}\bfseries\footnotesize~Sebastián~}}
\newcommand{\cj}[1]{\textcolor{juan}{\textbf{#1}}}
\newcommand{\cs}[1]{\textcolor{seb}{\textbf{#1}}}
\newcommand{\blk}[3]{%
  \par\medskip\noindent #1\ \textbf{#2}\hfill\textcolor{soft}{\small #3}\par
  \vspace{-4pt}\textcolor{soft!40}{\rule{\linewidth}{0.4pt}}\par\vspace{-2pt}}
\setlist[itemize]{leftmargin=1.4em,itemsep=1pt,topsep=2pt}
\pagestyle{plain}

\begin{document}

\begin{center}
  {\LARGE\bfseries Guion (20 min) --- El efecto Mpemba cuántico en HPC}\\[2pt]
  {\small\textcolor{soft}{Charla conjunta · versión extendida $\approx$20 min}}\\[6pt]
  {\footnotesize \Jb\ \textcolor{soft}{intro · T1 · redes tensoriales · conclusiones}\qquad
   \Sb\ \textcolor{soft}{evolución temporal · trayectorias · GPU · relevancia macro}}
\end{center}
\medskip\hrule\smallskip
{\footnotesize\textcolor{soft}{Lenguaje natural. Más detalle que la versión de 10 min:
incluye la intuición física, el razonamiento HPC y la lectura de cada figura.
Los tiempos suman $\approx$20:00.}}

% =============================================================== JUAN
\blk{\Jb}{0 · Apertura (\emph{Título})}{0:20}
\begin{itemize}
  \item ``Buenas. Vamos a contar un proyecto de \textbf{computación de alto
    rendimiento} aplicado a un fenómeno físico curioso: el \textbf{efecto Mpemba
    cuántico}. La idea: cuatro tareas numéricas, cada una con su versión \emph{serial}
    y \emph{paralela}, validadas contra soluciones exactas y con su estudio de escalado.''
  \item ``Nos repartimos: en \cj{teal} lo de Juan, en \cs{naranja} lo de Sebastián.''
\end{itemize}

\blk{\Jb}{1 · De lo clásico a lo cuántico --- La anomalía}{1:10}
\begin{itemize}
  \item ``El efecto Mpemba clásico es contraintuitivo: bajo ciertas condiciones, el
    agua \textbf{más caliente se congela antes} que la templada. Se observa también en
    coloides, granulares, imanes\dots''
  \item ``La clave \emph{no} es la temperatura en sí, sino una idea geométrica: si
    medimos una \textbf{distancia al equilibrio}, la curva que parte \textbf{más lejos}
    puede \textbf{cruzar} a la que parte más cerca y llegar antes. Ese \emph{cruce} es
    la firma del efecto.''
  \item ``Nuestra pregunta: ¿ocurre lo mismo en un sistema \textbf{cuántico abierto}?
    Y la respuesta es sí --- y además se puede \emph{predecir} con álgebra lineal.''
    (señalar la figura del cruce)
\end{itemize}

\blk{\Jb}{1 · Matemática: GKSL y autovalores}{1:40}
\begin{itemize}
  \item ``Un sistema cuántico abierto markoviano evoluciona con la \textbf{ecuación
    GKSL}, o de Lindblad: un término coherente (conmutador con $H$) más un
    \textbf{disipador} que modela el acoplamiento al baño.''
  \item ``Como es \emph{lineal} en $\rho$, la reescribimos con un operador, el
    \textbf{Liouvilliano} $\mathcal{L}$. Y aquí entra lo importante: diagonalizamos
    $\mathcal{L}$ y obtenemos sus \textbf{autovalores} $\lambda_k$, todos con parte
    real $\le 0$.''
  \item ``Entonces la evolución es una \textbf{suma de modos que decaen}: cada modo
    con su tasa $|\mathrm{Re}\,\lambda_k|$. A tiempos largos manda el \textbf{modo más
    lento}, $\lambda_2$.''
  \item ``El \textbf{criterio de Mpemba} es entonces puro álgebra: el peso inicial de
    cada modo es el solapamiento $a_k=\mathrm{Tr}(l_k^\dagger\rho_0)$. Si conseguimos
    una preparación con \cs{$a_2=0$}, esa preparación \emph{se salta el modo lento} y
    relaja a la tasa rápida --- por eso adelanta a otra que sí lo tiene.''
  \item ``Un detalle técnico clave: $\mathcal{L}$ es \textbf{no normal}, sus modos
    \emph{no} son ortogonales; eso es lo que permite que una curva que parte más lejos
    termine por debajo. Todo esto lo veremos reproducido numéricamente.''
  \item \emph{(Transición)} ``Lo primero, entonces: calcular esos autovalores. Tarea 1.''
\end{itemize}

\blk{\Jb}{1 · Predicción del cruce (\emph{sin evolucionar})}{0:50}
\begin{itemize}
  \item ``Antes de evolucionar nada, fijémonos en lo potente que es el espectro.
    Definimos dos preparaciones concretas, ambas relajando al \emph{mismo} estacionario:
    $|+\rangle^{\otimes N}$ --- todos los espines en $|+\rangle$, que parte \textbf{lejos}
    --- y $|0\rangle^{\otimes N}$, que parte \textbf{cerca}.''
  \item ``Con el espectro (los $\lambda_k$ y $\hat l_k$ que da T1) y los solapamientos
    $a_k=\mathrm{Tr}(\hat l_k^\dagger\rho_0)$ --- que son \emph{una simple traza}, sin
    integrar nada --- ya conocemos el destino de cada estado.''
  \item ``Y los números cantan: $|+\rangle$ arranca más alto ($D_0$: 0.99 vs 0.84) pero
    pesa \textbf{mucho menos} en el modo lento ($a_2$: 0.13 vs 0.50). Arranca alto y
    termina bajo $\Rightarrow$ \textbf{el cruce está garantizado}, con un $t^*$ predicho
    de $\approx0.68$.''
  \item ``Es decir: \cs{el espectro predice el Mpemba}; las tareas de evolución (T2, T3)
    solo lo van a \emph{confirmar} --- y son las que nos permiten llegar a muchos cuerpos
    donde ya no se puede diagonalizar.''
\end{itemize}

\blk{\Jb}{2 · T1 Modos lentos --- Método y discretización}{1:10}
\begin{itemize}
  \item ``Problema: el Liouvilliano es una matriz $4^N\times4^N$. Diagonalizarla en
    denso es inviable ya para pocos espines.''
  \item ``Solución: \textbf{Arnoldi}, pero con un truco. En vez de aplicar $\mathcal{L}$
    (cuyos modos de interés son los \emph{interiores}, cerca de 0, difíciles para
    Arnoldi), aplicamos el \textbf{propagador} $P=e^{\tau\mathcal{L}}$. Los modos lentos
    de $\mathcal{L}$ se vuelven los \textbf{dominantes} de $P$, que es justo lo que
    Arnoldi extrae primero --- lo llamamos \emph{faster than the clock}.''
  \item ``Y la acción de $P$ es \textbf{matrix-free}: integramos $\dot V=\mathcal{L}[V]$
    con RK4, sin formar nunca la matriz gigante.''
\end{itemize}

\blk{\Jb}{2 · T1 --- Núcleo serial\,$\vert$\,paralelo}{0:50}
\begin{itemize}
  \item ``El coste vive en el \textbf{producto de matrices} del SpMV. A la izquierda,
    serial: recorremos todas las filas.''
  \item ``A la derecha, paralelo en \textbf{dos niveles}: cada \textbf{rango MPI}
    calcula un \emph{bloque de filas} (paralelismo de datos), y dentro \textbf{OpenMP}
    reparte ese bloque entre hilos. Al final, \cs{MPI\_Allgatherv} reconstruye la
    matriz completa en todos los rangos.''
\end{itemize}

\blk{\Jb}{2 · T1 --- Resultados}{0:50}
\begin{itemize}
  \item ``Validación: los autovalores lentos coinciden con la diagonalización densa
    con errores de $10^{-13}$ a $10^{-5}$.''
  \item ``Escalado $\sim2.9\times$ con OpenMP y MPI. ¿Por qué no más? Porque el SpMV
    está \textbf{acoplado} y es \emph{memory-bound}: los hilos compiten por el ancho de
    banda, y MPI paga la comunicación del \texttt{Allgatherv}. Es el caso \emph{menos}
    paralelizable de los que veremos --- y es un resultado, no un fallo.''
  \item ``Y lo importante para la historia: T1 nos da el $\lambda_2$ y el $a_2$ que
    \textbf{determinan el cruce} de Mpemba.'' \emph{(pasa a Sebastián)}
\end{itemize}

% =============================================================== SEB
\blk{\Sb}{3 · T2a Evolución RK4 --- Método y discretización}{1:10}
\begin{itemize}
  \item ``Ahora la \textbf{evolución temporal}: dado un estado inicial, queremos la
    trayectoria $\rho(t)$ y su distancia al estacionario.''
  \item ``Vía directa: \textbf{Runge--Kutta de orden 4} sobre $\dot\rho=\mathcal{L}[\rho]$.
    Cuatro evaluaciones de $\mathcal{L}$ por paso, también \textbf{matrix-free}
    (conmutadores y disipador como productos de matrices $d\times d$).''
  \item ``¿Cómo paralelizar? Los pasos tienen \textbf{dependencia temporal estricta}:
    $\rho_{n+1}$ necesita $\rho_n$. No se puede paralelizar \emph{entre} pasos. El
    paralelismo está \emph{dentro} del paso, en el producto de matrices. Datos
    regulares en un nodo $\Rightarrow$ el paradigma natural es \textbf{OpenMP}, no MPI
    (que pagaría comunicación en cada uno de los miles de pasos).''
\end{itemize}

\blk{\Sb}{3 · T2a --- Núcleo serial\,$\vert$\,paralelo}{0:50}
\begin{itemize}
  \item ``El núcleo es el \texttt{matmul}. La versión paralela es \emph{una sola
    línea} de diferencia: \cs{\#pragma omp parallel for}, que reparte las filas $i$
    entre hilos. Son independientes, no hay condiciones de carrera.''
  \item ``Detalle bonito de ingeniería: es el \emph{mismo} binario; cambiando la
    variable \texttt{OMP\_NUM\_THREADS} obtenemos las curvas serie vs paralelo, sin
    recompilar.''
\end{itemize}

\blk{\Sb}{3 · T2a --- Resultados (computacional y físico)}{0:55}
\begin{itemize}
  \item ``Izquierda: tiempo serie vs paralelo según el tamaño. El speedup
    \textbf{crece con la malla} hasta $\sim4\times$, pero es \emph{sublineal}: otra vez
    \emph{bandwidth-bound} más el coste de crear hilos en cada \texttt{matmul}.''
  \item ``Derecha, lo físico: \textbf{aquí está el cruce}. La preparación que parte
    lejos (roja) \textbf{adelanta} a la cercana (azul) en \cs{$t^*\approx0.69$}.
    Validamos RK4 a orden 4 (pendiente $\approx4.04$).''
\end{itemize}

\blk{\Sb}{4 · T2b Krylov --- Método y el resultado central}{1:20}
\begin{itemize}
  \item ``Segunda vía para la misma evolución, más sofisticada: en vez de dar muchos
    pasitos, aplicamos directamente la \textbf{acción del exponencial}
    $e^{\tau\mathcal{L}}$, que es la solución exacta en un paso $\tau$.''
  \item ``Pero $e^{\tau\mathcal{L}}$ es un objeto enorme. El truco de \textbf{Krylov}:
    proyectamos en un subespacio pequeño (Arnoldi, $m\sim30$) y exponenciamos ahí una
    matriz $m\times m$ trivial.''
  \item ``Comparte el \emph{mismo} hotspot que RK4 (el \texttt{matmul}), así que la
    paralelización OpenMP y el escalado son equivalentes.''
  \item ``El resultado central está en esta gráfica: \textbf{a igual trabajo, Krylov es
    unas $25\times$ más preciso} que RK4, y permite pasos de tiempo grandes siendo
    \emph{incondicionalmente estable}. Es el método de elección para llegar al
    estacionario.''
\end{itemize}

\blk{\Sb}{4 · T2b --- El cruce}{0:35}
\begin{itemize}
  \item ``Y reproduce el \textbf{mismo cruce, $t^*\approx0.69$}, idéntico a RK4 --- de
    hecho coinciden a $10^{-10}$ en la distancia final. Mensaje: el cruce es
    \textbf{física}, no un artefacto del método numérico.''
\end{itemize}

\blk{\Sb}{5 · T3 Trayectorias --- Método y núcleo serial\,$\vert$\,paralelo}{1:10}
\begin{itemize}
  \item ``Hasta aquí trabajamos con $\rho$ completo, de tamaño $4^N$: inviable para
    \textbf{muchos cuerpos}. Cambio de paradigma: \textbf{trayectorias cuánticas}
    (Monte Carlo wavefunction).''
  \item ``En vez de propagar $\rho$, propagamos $M$ \textbf{estados puros}
    $|\psi\rangle$ de tamaño $2^N$ --- la raíz --- que evolucionan con un Hamiltoniano
    efectivo y dan \textbf{saltos} estocásticos. El promedio reconstruye $\rho$.''
  \item ``La gran ventaja para HPC: cada trayectoria es \textbf{independiente}
    --- \emph{embarrassingly parallel}. MPI reparte las $M$ entre rangos, cada uno
    calcula su lote \emph{sin comunicación}, y \textbf{una sola \cs{MPI\_Reduce}} al
    final promedia.''
\end{itemize}

\blk{\Sb}{5 · T3 --- Resultados}{0:50}
\begin{itemize}
  \item ``Y se nota: el escalado es \textbf{casi ideal, $\sim5.6\times$ y constante}
    al crecer el trabajo --- en claro contraste con el OpenMP sublineal de T2. Ese
    contraste es \emph{el mensaje HPC de la charla}: \textbf{el patrón de cómputo dicta
    el paradigma}, no una preferencia.''
  \item ``La precisión converge como $M^{-1/2}$. Y --- otra vez --- el \textbf{cruce de
    Mpemba}, ahora con el método estocástico, en \cs{$t^*\approx0.75$} (la pequeña
    diferencia es el ruido Monte Carlo).''
\end{itemize}

\blk{\Sb}{5 · CUDA T4 --- GPU}{0:55}
\begin{itemize}
  \item ``Como bonus, el \textbf{tercer escalón} de paralelización para la evolución
    densa: la GPU. El producto de matrices denso es el caso ideal para cuBLAS.''
  \item ``En una \textbf{T4} de Colab: pierde para tamaños pequeños (overhead de lanzar
    kernels) pero \textbf{gana $\sim5\times$} en cuanto la malla crece --- para $N=9$,
    de 943 segundos en CPU a 190 en GPU. El producto puro llega a 250 GFLOP/s.''
  \item ``Y validación triple: GPU $=$ CPU $=$ C++. El notebook ejecutado y los
    resultados están en el repositorio.'' \emph{(pasa a Juan)}
\end{itemize}

\blk{\Jb}{6 · T3b Redes tensoriales --- Método y discretización}{1:10}
\begin{itemize}
  \item ``La otra gran ruta a muchos cuerpos, complementaria a las trayectorias:
    \textbf{redes tensoriales}.''
  \item ``Escribimos la matriz densidad como un \textbf{superket} en forma de MPS ---
    una cadena de tensores, con dimensión física 4 por sitio. La ecuación maestra se
    integra con \textbf{TEBD}: trotterizamos $e^{dt\,\mathcal{L}}$ en \textbf{compuertas
    locales} de dos sitios y \textbf{truncamos} el enlace a $\chi$ con una SVD.''
  \item ``Resultado: coste \textbf{polinómico en $\chi$}, no exponencial en $N$. Eso
    abre tamaños imposibles en denso.''
  \item ``Paralelización por \textbf{hilos}: en una capa de Trotter, los bonds pares
    (o impares) son \emph{disjuntos}; sus SVD son independientes. Y numpy
    \textbf{libera el GIL} durante la SVD, así que los hilos escalan de verdad.''
\end{itemize}

\blk{\Jb}{6 · T3b --- Núcleo serial\,$\vert$\,paralelo}{0:45}
\begin{itemize}
  \item ``Serial: un bucle que aplica las compuertas de los bonds una a una. Paralelo:
    un \cs{executor.map} lanza las compuertas disjuntas en hilos.''
  \item ``Es un paralelismo \textbf{estructurado} --- ni tan acoplado como el SpMV de
    T1, ni tan libre como las trayectorias. Queda en medio, y eso se ve en el escalado.''
\end{itemize}

\blk{\Jb}{6 · T3b --- Resultados}{0:50}
\begin{itemize}
  \item ``Validación: comparado con el resultado exacto a $N$ pequeño, el error
    \textbf{cae al subir $\chi$}, hasta $6\times10^{-7}$ --- control sistemático.''
  \item ``Escala $\sim3.6\times$ en hilos (entre T1 y las trayectorias, como
    anticipamos). Y el alcance: llegamos a \textbf{$N=32$} --- un espacio de Hilbert de
    $10^{19}$, absolutamente imposible en denso --- evolucionado en segundos.''
\end{itemize}

\blk{\Jb}{6 · T3b --- El cruce de Mpemba}{0:45}
\begin{itemize}
  \item ``Y cerramos con lo que une todo: \textbf{también las redes tensoriales
    reproducen el cruce}. Las mismas dos preparaciones, $|+\rangle$ que parte lejos
    adelanta a $|0\rangle$ en \cs{$t^*\approx0.69$}.''
  \item ``Idéntico a RK4 y Krylov, porque a $\chi$ grande TEBD es exacto. Es decir:
    \textbf{los cuatro métodos de evolución} --- RK4, Krylov, trayectorias y redes
    tensoriales --- dan \emph{el mismo} cruce. Validación cruzada total.''
\end{itemize}

\blk{\Jb}{7 · Conclusiones generales}{1:20}
\begin{itemize}
  \item ``\textbf{HPC.} El mensaje central: el \emph{patrón de cómputo} determina la
    escalabilidad. Lo vimos como un espectro:
    \emph{acoplado} (T1, SpMV, $\sim3\times$) $<$ \emph{estructurado} (T3b, redes
    tensoriales, $\sim3.6\times$) $<$ \emph{independiente} (trayectorias, $\sim5.6\times$
    casi ideal). Y la evolución densa, además, vuela en GPU.''
  \item ``Cada paradigma --- OpenMP, MPI, hilos, CUDA --- lo elegimos por la
    \emph{estructura del algoritmo}, no por moda.''
  \item ``\textbf{Física.} El efecto Mpemba cuántico tiene origen
    \textbf{geométrico-espectral}: la no-normalidad de $\mathcal{L}$ más el criterio
    $a_2=0$. Lo \textbf{reprodujimos con los cuatro métodos} ($t^*\approx0.69$--$0.75$),
    validados de forma cruzada, y escalado a muchos cuerpos.'' \emph{(pasa a Sebastián)}
\end{itemize}

\blk{\Sb}{8 · Relevancia en el efecto macro/meso}{1:00}
\begin{itemize}
  \item ``Para cerrar, por qué importa. La \textbf{firma es universal}: el mismo
    mecanismo de \emph{geometría espectral de la relajación} explica el cruce desde el
    macro clásico (agua, coloides) hasta lo cuántico. Es el mismo lenguaje matemático.''
  \item ``Lo que devuelve el enfoque cuántico al problema clásico: un \textbf{criterio
    operacional} --- $a_2=0$ --- y las \textbf{herramientas HPC} para \emph{predecir} y
    \emph{diseñar} el efecto en vez de buscarlo a ciegas.''
  \item ``Con aplicaciones concretas: \textbf{termometría cuántica} y protocolos de
    \textbf{enfriamiento acelerado}.''
  \item ``Y esa es la idea: una anomalía vieja de la física, entendida como álgebra de
    autovalores y llevada a HPC. Muchas gracias.'' \emph{(cierre)}
\end{itemize}

\medskip\hrule\smallskip
\noindent\textbf{\footnotesize Notas.}
\begin{itemize}[font=\footnotesize]\footnotesize
  \item Énfasis recurrente: los \textbf{cuatro cruces} (T2a, T2b, T3, T3b dan el mismo
    $t^*$) y el contraste de escalado \textbf{acoplado $<$ estructurado $<$ independiente}.
  \item Si vas sobre tiempo: la slide CUDA es la más prescindible; las ecuaciones de la
    slide matemática se pueden resumir.
  \item Frase ancla: ``\textbf{el patrón de cómputo dicta el paradigma}''.
\end{itemize}

\end{document}
